\paragraph{Next order at \(N_s=6\), \(n_D=2\): exact \(\mathfrak J=0\) backgrounds and the surviving quartic (conditional; not closure).}
Keep the setting of the paragraphs on the gauge parameter and on fermionic content: Poisson limit, Euclidean section \(\eta=\delta\), \(\alpha=1\), one loop, \(N_s=6\) scalars, the connection with its ghosts and \(n_D=2\) Dirac fields, so that \(N_{\mathrm{tot}}=0\) and the weighted coefficient \(\kappa_X\) of \(\rho^2X\) vanishes. The statements below are local, on regions where \(\operatorname{Pf}\omega\ne0\). Write \((b_\sigma)_{\mu\nu}=G_{\sigma\tau}\Theta^{\rho\tau}\partial_\rho\Theta_{\mu\nu}\), the \(\mathfrak{so}(4)\)-valued one-form with \(\omega_\sigma=-\tfrac12\partial_\sigma\ln\rho\,\mathbb I+b_\sigma\) in the vector sector and, through the spinor lift \(\beta_\sigma=\tfrac14\gamma^{\mu\nu}(b_\sigma)_{\mu\nu}\), in the Dirac sector. For every Poisson bivector \(\mathfrak J^\nu=\Box_Gx^\nu-G^{\mu\nu}\partial_\mu\ln\rho\) (exact check at random two-jets), so that \(\mathfrak J=0\) says that the coordinates are harmonic for \(G\) up to the gradient of \(\ln\rho\).

\emph{Backgrounds with \(\mathfrak J=0\).} Up to a constant factor, \(\rho\,\Theta^{\mu\nu}\Theta^{\mu\nu}\) is the Poisson-limit density of the first term of Equation~\eqref{eq:Smatrix}, and \(\rho\,\Theta^{\mu\nu}\Theta^{\mu\nu}\ge4\) pointwise, with equality if and only if \(\Theta\) is self-dual or anti-self-dual. For such \(\Theta\), \(\Theta\Theta^T\) is a multiple of the identity, so \(G^{\mu\nu}=\rho^{-1}\delta^{\mu\nu}\), \(g_\rho=\delta\) and \(\mathfrak J=0\) by the identity above. A self-dual \(\Theta\) is the inverse of a closed self-dual two-form \(\omega=b_i\eta^i\), with \(\eta^i_{ab}=\epsilon_{iab}\), \(\eta^i_{a4}=\delta_{ia}\); closure is \(\operatorname{div}b=0\), \(\partial_4b=-\operatorname{curl}b\), and then \(\rho=|b|^2\), \(G=\rho\,\delta\) is conformally flat with \(R_G\ne0\) in general, while \(\rho R[g_\rho]=0\). Since \(\partial_\rho\Theta^{\mu\nu}\) is then self-dual in \((\mu\nu)\), \(b_\sigma\) takes values in one \(\mathfrak{su}(2)\) factor of \(\mathfrak{so}(4)\); the same holds for anti-self-dual \(\Theta\). An example invariant under the rotations of the \((x^1,x^2)\) plane, which fix the plane \(\{x^1=x^2=0\}\) pointwise, is \(b_1+\mathrm ib_2=-(x^1+\mathrm ix^2)(x^3-\mathrm ix^4)\), \(b_3=1+(x^3)^2-(x^4)^2\) (exact check of closure, invariance, \(\mathfrak J=0\) and \(R_G\ne0\) at sample points; the vector \((b_1,b_2,b_3)\) vanishes on the locus \(x^1=x^2=0\), \((x^4)^2=1+(x^3)^2\), where \(\operatorname{Pf}\omega=0\) and the example is not defined).

Outside this class, let \(\Lambda^i=\eta^i_{\mu\nu}x^\mu\mathrm dx^\nu\), so that \(\Lambda^3=x^1\mathrm dx^2-x^2\mathrm dx^1+x^3\mathrm dx^4-x^4\mathrm dx^3\), and take \(\omega=\mathrm d(h(r)\Lambda^3)\), invariant under the rotations of both coordinate planes. Then \(\mathfrak J\) has only a radial component, proportional to \(2h^2rh''+10h^2h'+4hrh'^2+r^2h'^3\), and \(u=rh'/h\) obeys \(r\,\mathrm du/\mathrm dr=-\tfrac12u(u+2)(u+4)\). The fixed points \(u=0\) and \(u=-4\) are the self-dual and anti-self-dual solutions, \(u=-2\) is degenerate (\(\operatorname{Pf}\omega\propto h^2(u+2)\)), and the other solutions are \(h=c\,(1\pm\sqrt{1-C/r^4})\), \(C\in\mathbb R\) (checked at generic Cartesian points), locally, for \(r^4>C\) when \(C>0\); at \(r^4=C\) the metric \(G\) degenerates. On them \(G=\beta^2[s^{-2}(\mathrm dr^2+e_3^2)+e_1^2+e_2^2]\) with \(\beta=2h\), \(s^2=1-C/r^4\), \(e_i=\Lambda^i/r\), and \(C_{\mu\nu\rho\sigma}C^{\mu\nu\rho\sigma}=\tfrac32u^2(u+4)^2/(hr)^4\), which vanishes only at the two end points.

\emph{Coefficients.} Let \(\mathcal S_4=\operatorname{tr}_V[(b_\sigma b^\sigma)^2]-n_D\operatorname{tr}_D[(\beta_\sigma\beta^\sigma)^2]\), with \(b_\sigma b^\sigma=G^{\sigma\tau}b_\sigma b_\tau\) and traces over the four vector and the four Dirac components. In the variables \((G,\rho,b)\) the endomorphism of every sector is \(E=e\,\mathbb I+X+b_\sigma b^\sigma\), with the same scalar \(e=-\tfrac12\Box_G\ln\rho+\tfrac14|\partial\ln\rho|^2_G\), \(X\) linear in \(b\) and traceless, and the bundle curvature is that of \(b\). The weighted sums of \(\operatorname{tr}\mathbb I\) and of the Dynkin index of \(b\) are \(N_{\mathrm{tot}}=0\) and \(\kappa_X=0\), and \(\mathfrak{so}(4)\) has no cubic invariant. Hence, for every Poisson bivector, the weighted sum of the coefficients \(a_1\) vanishes pointwise and the weighted sum of the coefficients \(a_2\) is \(\tfrac12\mathcal S_4\), up to total derivatives, which vanish by the same sums: no \(R^2\), \(R_{\mu\nu}R^{\mu\nu}\), Riemann-squared, Weyl-squared or Euler term, no \(RE\) term and no term of Einstein--Hilbert type survives (exact rational arithmetic at random two-jets, also for \(\eta=\operatorname{diag}(-1,1,1,1)\); on the self-dual and anti-self-dual backgrounds the weighted sum of \(a_2\) vanishes while the single-sector values do not, and it does so for the backgrounds tested only at \((N_s,n_D)=(6,2)\), with the vector weight \(1\) and the ghost weight \(-2\) fixed in that scan). With \(l^k_\sigma=\tfrac14\eta^k_{\mu\nu}(b_\sigma)_{\mu\nu}\) and \(r^k_\sigma=\tfrac14\bar\eta^k_{\mu\nu}(b_\sigma)_{\mu\nu}\) the self-dual and anti-self-dual parts (\(\bar\eta^k_{ab}=\epsilon_{kab}\), \(\bar\eta^k_{a4}=-\delta_{ka}\)), and \(\mathsf A=G^{\sigma\tau}l_\sigma l_\tau^T\), \(\mathsf B=G^{\sigma\tau}r_\sigma r_\tau^T\), \(\mathsf C=G^{\sigma\tau}l_\sigma r_\tau^T\),
\[
  \mathcal S_4=8\operatorname{tr}\mathsf A\,\operatorname{tr}\mathsf B+16\operatorname{tr}(\mathsf C\mathsf C^T)\ \ge0,
\]
with equality at a point if and only if \(b_\sigma\) takes values in one \(\mathfrak{su}(2)\) factor there. With one common proper-time limit for all sectors, the coefficient of \(\ln\LUV\) in the induced action at this content is therefore \(\propto\int\rho^2\mathcal S_4\), quartic in \(\partial\Theta\); finite terms of the same order are not controlled. It vanishes on the self-dual and anti-self-dual backgrounds above. On the \(\mathrm d(h\Lambda^3)\) family \(b_\sigma\) was found to be anti-self-dual-valued in every sampled case with \(\mathfrak J=0\), so \(\mathcal S_4=0\) there, as the weighted \(a_2\) computed there confirms (not proved), and it has both parts in the sampled cases with \(\mathfrak J\ne0\). The condition \(\mathfrak J=0\) does not by itself force \(\mathcal S_4=0\): about a constant \(\bar\Theta\), with \(\Theta=(\psi^*\bar\omega)^{-1}\) and \(\psi=x+\epsilon a\), the condition \(\mathfrak J=O(\epsilon^2)\) is the source-free Maxwell equation \(D^\mu(D^\mu a^\nu-D^\nu a^\mu)=0\), \(D^\mu=\bar\Theta^{\mu\rho}\partial_\rho\); for quadratic \(a\) its solutions give constant \(b_\sigma=O(\epsilon)\) with both parts and \(\mathcal S_4>0\) (random solutions for one \(\bar\Theta\); in one example the order \(\epsilon^2\) of \(\mathfrak J=0\) is solved by a cubic correction of \(\psi\)). The convergence of the series, a global background, the domains of the intertwining lemma above and the tip domain are not examined. On backgrounds with \(\mathfrak J=0\) globally the logarithmic coefficient does not depend on \(\alpha\), because the longitudinal sector is a scalar tower with proper time rescaled and the coefficient of \(s^2\) is unchanged.

\emph{Where an Einstein--Hilbert term can enter.} In the massless class every local term of order \(\LUV^0\) has four derivatives. If a sector carries a constant mass, \(\mathcal O_i\to\mathcal O_i+m_i^2\) (a formal shift of the operator: no gauge-invariant mass term and no ghost or longitudinal consistency is asserted), then \(E_i\to E_i+m_i^2\), and the coefficient \(a_2\) of that sector changes by \(-m_i^2a_1^{(i)}+\tfrac12\operatorname{tr}\mathbb I\,m_i^4\) (exact check). The curvature part of the change of the weighted sum is \(-\bigl(\sum_iw_in_im_i^2/6\bigr)\rho R[g_\rho]-\bigl(\sum_iw_i\kappa_im_i^2\bigr)\rho^2X\), where \(n_i\) is the number of components of the sector and \(\kappa_i\) its weight of \(\rho^2X\) (\(1\) for the connection, \(\tfrac12\) per Dirac field, \(0\) otherwise). It vanishes for one common mass, because \(N_{\mathrm{tot}}=0\) and \(\kappa_X=0\), and is nonzero only if bosonic and fermionic masses differ. The inert transverse matrices of this class carry no such scale. In the construction of~\cite{Steinacker2023} the Einstein--Hilbert term is generated by a mixed quartic term with a sum over the Kaluza--Klein modes of fuzzy extra dimensions, and the Newton constant is set by their scale; the class above has no analogue of it.

These statements are for the Poisson-limit operators, the Euclidean section, one loop, \(\alpha=1\), the content \(N_s=6\), \(n_D=2\) with the Yukawa couplings dropped, and the backgrounds named. They compute no logarithmic coefficient of the cone (area) side, so that no comparison of \(G_T\) with \(G_A\) at logarithmic order follows; they do not derive the content, the state, the tip domain or Assumption~\ref{ass:emergent-op}, do not treat Moyal corrections or fuzzy extra dimensions, and give no value or sign of \(G_{\mathrm{eff}}\). Problems~B, C, D, F and~G remain open.
